\section{ZPPO algorithm}
\label{sec:appendix:algorithm}

Algorithm~\ref{alg:zppo} summarizes one ZPPO training step. The key invariant is that every response used in the policy gradient is generated by the current student. The frozen teacher is queried only to provide prompt-side candidate text for BCQ (and to compress candidates under a shared compression prompt and token cap), so teacher-generated tokens never enter the gradient as response tokens.

ZPPO inherits its RL backbone from three lines of recent work and adds four novel components. \emph{Inherited from GRPO~\citep{shao2024deepseekmath}}: (G1)~group-relative advantage formulation; (G2)~PPO-style clipped surrogate objective. \emph{Inherited from DAPO~\citep{yu2025dapo}}: (Da1)~asymmetric clip-higher with $(\epsilon_{\rm low},\epsilon_{\rm high})\!=\!(0.20, 0.28)$; (Da2)~token-level loss aggregation; (Da3)~no KL penalty against a reference policy. \emph{Inherited from REINFORCE++~\citep{hu2025reinforce++}}: (R1)~two-step advantage estimator (Eqs.~\ref{eq:adv-step1}--\ref{eq:adv-step2}). \emph{Our recipe choice on top of REINFORCE++}: (Z0)~zero-advantage-group exclusion from the batch-statistics computation in Step~2, ablated in Sec.~\ref{sec:experiments:recipe}(ii) and derived in closed form in Eqs.~\ref{eq:groupsum-zero}--\ref{eq:wzero-amplify}. \emph{ZPPO's own contributions}: (Z1)~BCQ -- the anonymized binary-candidate prompt reformulation that turns a hard question into a discriminative one-of-two judgment between a correct teacher trace and a wrong student trace, both compressed by the same teacher under a shared compression prompt and token cap to mitigate obvious surface cues in the discrimination; (Z2)~NCQ -- the collective-negative-candidate prompt reformulation that surfaces \emph{all} of the student's wrong rollouts on the same question, together with their parsed final answers, so the student must avoid the wrong rollouts it has just produced; (Z3)~prompt replay buffer with hard-question admission, graduation by mean rollout accuracy, and FIFO eviction; (Z4)~the super-additive combination of (Z3) with (Z1)+(Z2) that is the central empirical claim of the paper (Sec.~\ref{sec:experiments:ablation}). Algorithm~\ref{alg:zppo} annotates each step with the tags above.

\begin{algorithm*}[!htbp]
\SetAlgoLined
\DontPrintSemicolon
\KwIn{Student $\pi_\theta$, teacher $\pi_{\rm T}$, dataset $\mathcal{D}$, buffer $\mathcal{B}$; $G_{\rm S}, G_{\rm T}, \rho_{\rm replay}, \rho_{\rm aug}, |\mathcal{B}|_{\max}, I, \tau\!=\!0.5$.}
\KwOut{Updated $\pi_\theta$, $\mathcal{B}$.}

\textbf{// 1. Plain rollouts.} \textit{[G1, Z3]}\;
$X\!\leftarrow\!X_{\rm new}\!\cup\!X_{\rm replay}$ with $X_{\rm new}\!\sim\!\mathcal{D}$, $X_{\rm replay}\!\sim\!\mathcal{B}$, $|X_{\rm replay}|\!=\!\rho_{\rm replay}|X_{\rm new}|$. For each $x\!\in\!X$, draw $\{y_{\rm S}^{(g)}\}_{g=1}^{G_{\rm S}}\!\sim\!\pi_\theta(\cdot|x)$, score, compute $\bar{r}_x$, and collect wrong subset $\{y_{\rm S}^{(-)}(x)\}$.\;

\textbf{// 2. Teacher rollouts on hard questions.} \textit{[Z1]}\;
$X_{\rm hard}\!\leftarrow\!\{x:\bar{r}_x\!<\!\tau\}$. For each $x\!\in\!X_{\rm hard}$, draw $\{y_{\rm T}^{(g)}(x)\}_{g=1}^{G_{\rm T}}\!\sim\!\pi_{\rm T}(\cdot|x)$, score, and keep the correct subset $\{y_{\rm T}^{(+)}(x)\}$.\;

\textbf{// 3. Pre-cap on base questions.} \textit{[Z1, Z2]}\;
Keep the top $\rho_{\rm aug}|X_{\rm new}|$ of $X_{\rm hard}$ by ascending $\bar{r}_x$ as $X_{\rm aug}^{\rm pre}$ (every hard $x$ has $\{y_{\rm S}^{(-)}\}\!\neq\!\emptyset$ by $\bar{r}_x\!<\!\tau\!=\!0.5$; the per-instance BCQ/NCQ admissibility check happens in Step~4).\;

\textbf{// 4. Build BCQ/NCQ prompts.} \textit{[Z1, Z2]}\;
$\mathcal{A}\!\leftarrow\!\emptyset$. \ForEach{$x \in X_{\rm aug}^{\rm pre}$}{
    \lIf{$\{y_{\rm T}^{(+)}\},\{y_{\rm S}^{(-)}\}\!\neq\!\emptyset$}{teacher-compress one $y_{\rm T}^{(+)}$ and one $y_{\rm S}^{(-)}$ under the shared compression prompt and token cap, shuffle, form $x_{\rm BCQ}$; add $(x_{\rm BCQ},\mathrm{uid}_{\rm BCQ},x)$ to $\mathcal{A}$}
    \lIf{$\{y_{\rm S}^{(-)}\}\!\neq\!\emptyset$}{teacher-rewrite all wrong rollouts, list their parsed answers, form $x_{\rm NCQ}$; add $(x_{\rm NCQ},\mathrm{uid}_{\rm NCQ},x)$ to $\mathcal{A}$}
}

\textbf{// 5. Post-cap and reformulated rollouts.} \textit{[Z1, Z2]}\;
\If{$|\mathcal{A}|\!>\!\rho_{\rm aug}|X_{\rm new}|$}{
   Build $\mathcal{A}'\!\leftarrow\!\emptyset$ by iterating $x\!\in\!X_{\rm aug}^{\rm pre}$ in ascending $\bar{r}_x$; for each $x$, append its BCQ instance (if present) then its NCQ instance (if present) to $\mathcal{A}'$. If appending the next instance would push $|\mathcal{A}'|$ above $\rho_{\rm aug}|X_{\rm new}|$, skip that single instance (so when only one slot remains in a question that has both, BCQ is kept and NCQ is dropped) and stop. Set $\mathcal{A}\!\leftarrow\!\mathcal{A}'$.\;
}
For each $(x',\mathrm{uid}',x)\!\in\!\mathcal{A}$, draw $\{y^{(g)}\}_{g=1}^{G_{\rm S}}\!\sim\!\pi_\theta(\cdot|x')$ under $\mathrm{uid}'$ and score.\;

\textbf{// 6. Update and buffer refresh.} \textit{[G2, Da1--Da3, R1, Z0, Z3, Z4]}\;
Treat plain, BCQ, NCQ rollouts on $x$ as three separate groups (group key $=$ uid), each of size $G_{\rm S}$. Apply Eqs.~\eqref{eq:adv-step1}--\eqref{eq:adv-step2} (Step~1: subtract group mean; Step~2: batch-normalize \emph{only over the non-trivial subset} $\mathcal{G}^{\star}\!=\!\{g\!:\!\mathrm{std}_x\!>\!0\}$, leaving trivial groups at $A^{(g)}\!=\!0$); update $\pi_\theta$ for $I$ iterations.\;
Compute $\bar{r}_x$ from \emph{plain} rollouts only; admit $\{x:\bar{r}_x\!<\!\tau\}$ to $\mathcal{B}$, graduate the rest, FIFO-evict until $|\mathcal{B}|\!\leq\!|\mathcal{B}|_{\max}$.\;

\Return{$\pi_\theta$, $\mathcal{B}$}.
\caption{ZPPO training step. Tags: G1--G2 from GRPO~\citep{shao2024deepseekmath}, Da1--Da3 from DAPO~\citep{yu2025dapo}, R1 from REINFORCE++~\citep{hu2025reinforce++}; Z0 marks our recipe choice on top of REINFORCE++ (zero-advantage-group exclusion); Z1--Z4 are ZPPO's contributions (BCQ, NCQ, prompt replay buffer, and their super-additive combination).}
\label{alg:zppo}
\end{algorithm*}

A few practical notes are worth pointing out. First, BCQ requires \emph{both} a correct teacher rollout and a wrong student rollout to construct a candidate pair; questions on which the teacher itself fails on every $y_{\rm T}$ skip BCQ for that visit and only contribute NCQ (or, if all student rollouts succeed on the second try, neither). Second, candidate compression runs in parallel with Stage~1 scoring, so it does not stall the rollout loop. Third, the teacher rollouts $\{y_{\rm T}\}$ on a replayed question are re-sampled on every visit -- the BCQ candidate seen by the student on visit $k$ is not the same as the one seen on visit $k\!-\!1$, even when the underlying question is identical. This freshness is what allows the buffer to avoid degenerating into a single fixed exemplar per hard question. Fourth, only the \emph{plain} student rollouts on the original $x$ feed buffer admission and graduation in Stage~6: the BCQ/NCQ rollouts contribute to the gradient but are skipped when computing the per-question accuracy used to decide buffer membership, so a question is judged ``mastered'' only when the student can solve it \emph{without} candidate references. Fifth, the three groups (plain, BCQ, NCQ) on the same hard $x$ each have their own group identifier and are advantage-normalized independently; the plain group on $x$ and the reformulated groups on $x_{\rm BCQ}$, $x_{\rm NCQ}$ are sampled from different prompt distributions and never share a baseline.

\paragraph{What counts as a ``rollout step,'' a ``gradient update,'' and a ``token.''} Because $\rho_{\rm replay}$ and $\rho_{\rm aug}$ are both expressed as a fraction of \texttt{rollout\_batch\_size}, the prompt counts seen by the gradient at each rollout step are additive rather than fractions of a fixed budget. One rollout step contains: (i)~$X_{\rm new}$ new prompts drawn from $\mathcal{D}$ (over the rollout-step budget in Tab.~\ref{tab:appendix:hyperparameters}, $X_{\rm new}$ alone consumes the entire ZPPO-$77$K corpus in a single pass; $X_{\rm replay}$ does not count against the dataset budget); (ii)~$X_{\rm replay}$ prompts drawn from $\mathcal{B}$ on top of $X_{\rm new}$ with $|X_{\rm replay}|\!=\!\rho_{\rm replay}|X_{\rm new}|$; and (iii)~at most $\rho_{\rm aug}|X_{\rm new}|$ BCQ/NCQ instances after the combined post-cap of Step~5. The cap is enforced in two stages. (a)~\emph{Pre-cap on base questions} (Step~3 of Alg.~\ref{alg:zppo}): once teacher rollouts have been drawn on $X_{\rm hard}$ in Step~2, we keep the top $\rho_{\rm aug}|X_{\rm new}|$ base questions by ascending $\bar{r}_x$ that admit at least one of the BCQ or NCQ branches. (b)~\emph{Post-cap on combined BCQ$+$NCQ instances} (Step~5 of Alg.~\ref{alg:zppo}): because each $x\!\in\!X_{\rm aug}^{\rm pre}$ may contribute one BCQ instance (when $\{y_{\rm T}^{(+)}\}\!\neq\!\emptyset$ and $\{y_{\rm S}^{(-)}\}\!\neq\!\emptyset$) and one NCQ instance (when $\{y_{\rm S}^{(-)}\}\!\neq\!\emptyset$), the pre-instance count $|\mathcal{A}|$ can be up to twice the pre-cap; the \texttt{aug\_max\_ratio} clause in the code then prunes $\mathcal{A}$ as a single combined BCQ$+$NCQ pool (not separately per type) down to at most $\rho_{\rm aug}|X_{\rm new}|$. Each surviving prompt -- plain new, replayed, BCQ, or NCQ -- carries $G_{\rm S}$ rollouts, and the actual per-step prompt count enters the gradient via gradient accumulation; FLOPs in Tab.~\ref{tab:appendix:compute} are accumulated over \emph{all} gradient-counted tokens (plain $+$ replay $+$ BCQ $+$ NCQ, with the BCQ$+$NCQ portion respecting the augmentation cap above), with the $I$-update multiplier applied to each.

\paragraph{Two-step advantage estimator (REINFORCE++) used by ZPPO.}
Eq.~\eqref{eq:grpo-adv} in Sec.~\ref{sec:method:prelim} is the textbook GRPO advantage that ZPPO builds on. We adopt the two-step variant of \citet{hu2025reinforce++} (their \emph{REINFORCE++}, Eqs.~6--7), which decouples within-group centering from across-group normalization; for clarity, we restate it in our notation and make explicit which sub-population enters the batch statistics in Step~2. Let $\mathcal{G}$ denote the set of groups in a mini-batch (each group $g$ has $G_{\rm S}$ rollouts sharing a uid), and let $\mathcal{G}^{\star}\!\subseteq\!\mathcal{G}$ be the \emph{non-trivial} subset, defined as those groups with $\mathrm{std}_x\!>\!0$ (i.e.\ neither all-correct nor all-wrong). Step~1 subtracts the group mean,
\begin{equation}
A'^{(g)}_{x,i} \;=\; r(x, y^{(g,i)}) - \bar{r}_x,
\label{eq:adv-step1}
\end{equation}
which produces $A'\!=\!0$ for every rollout in a trivial (zero-advantage) group by construction. Step~2 batch-normalizes across a sub-population $\mathcal{S}\!\subseteq\!\mathcal{G}$ of groups:
\begin{align}
A^{(g)}_{x,i} &\;=\; \frac{A'^{(g)}_{x,i} - \mu_{\mathcal{S}}}{\sigma_{\mathcal{S}} + \epsilon}, \nonumber \\
\mu_{\mathcal{S}} &\;=\; \mathrm{mean}_{(g,i):g\in\mathcal{S}}\,A'^{(g)}_{x,i}, \nonumber \\
\sigma_{\mathcal{S}} &\;=\; \mathrm{std}_{(g,i):g\in\mathcal{S}}\,A'^{(g)}_{x,i}.
\label{eq:adv-step2}
\end{align}
The three settings compared in Sec.~\ref{sec:experiments:recipe}(ii) (\emph{No norm}, \emph{Norm w/o Zero}, \emph{Norm w/ Zero}) all perform Step~1 and differ only in $\mathcal{S}$ and in whether Step~2 is applied:
\begin{itemize}[leftmargin=12pt,topsep=2pt,itemsep=2pt]
  \item \emph{No norm}: Step~2 is skipped entirely, $A^{(g)}\!=\!A'^{(g)}$ (within-group centering only); equivalent to the standard group-relative advantage of Eq.~\eqref{eq:grpo-adv} up to the $\mathrm{std}_x$ rescaling.
  \item \emph{Norm w/o Zero} (ZPPO default): $\mathcal{S}\!=\!\mathcal{G}^{\star}$. The batch statistics in Step~2 are computed over the non-trivial groups only, and trivial groups keep $A^{(g)}\!=\!0$ (since their $A'\!=\!0$ is left untouched).
  \item \emph{Norm w/ Zero}: $\mathcal{S}\!=\!\mathcal{G}$. The batch statistics in Step~2 include every group, even the trivial ones whose $A'\!=\!0$. We make the consequence quantitative. Because Step~1 centers each group to a sum of zero,
  \begin{equation}
  \sum_{i=1}^{G_{\rm S}} A'^{(g,i)} \;=\; \sum_{i=1}^{G_{\rm S}} \bigl(r^{(g,i)}-\bar{r}_g\bigr) \;=\; 0 \quad \text{for every } g\in\mathcal{G},
  \label{eq:groupsum-zero}
  \end{equation}
  the batch mean is identically zero, $\mu_{\mathcal{G}}\!\equiv\!0$, irrespective of how many trivial groups the batch happens to contain. The batch standard deviation, on the other hand, is depressed by the trivial groups' zeros. Writing $f_{\rm nt}\!=\!|\mathcal{G}^{\star}|/|\mathcal{G}|$ and using equal group size $G_{\rm S}$,
  \begin{align}
  \sigma_{\mathcal{G}}^{2}
    &\;=\; \frac{1}{|\mathcal{G}|\,G_{\rm S}} \sum_{g\in\mathcal{G}}\sum_{i} \bigl(A'^{(g,i)}\bigr)^{2} \nonumber \\
    &\;=\; \frac{|\mathcal{G}^{\star}|}{|\mathcal{G}|}\cdot\frac{1}{|\mathcal{G}^{\star}|\,G_{\rm S}}\sum_{g\in\mathcal{G}^{\star}}\sum_{i} \bigl(A'^{(g,i)}\bigr)^{2} \nonumber \\
    &\;=\; f_{\rm nt}\,\sigma_{\mathcal{G}^{\star}}^{2},
  \label{eq:sigma-shrink}
  \end{align}
  so $\sigma_{\mathcal{G}}=\sqrt{f_{\rm nt}}\,\sigma_{\mathcal{G}^{\star}}$. Substituting into Eq.~\eqref{eq:adv-step2} gives a closed form for the resulting advantages:
  \begin{align}
  A^{(g,i)}_{\text{w/Zero}}
    &\;=\; \frac{r^{(g,i)}-\bar{r}_g}{\sqrt{f_{\rm nt}}\,\sigma_{\mathcal{G}^{\star}}+\epsilon} \nonumber \\
    &\;=\; \frac{1}{\sqrt{f_{\rm nt}}}\cdot A^{(g,i)}_{\text{ZPPO}} \quad (g\in\mathcal{G}^{\star}), \nonumber \\
  A^{(g,i)}_{\text{w/Zero}}
    &\;=\; 0 \quad (g\notin\mathcal{G}^{\star}).
  \label{eq:wzero-amplify}
  \end{align}
  Two things follow. First, trivial groups stay at $A^{(g,i)}\!=\!0$ exactly as under ZPPO -- they receive neither a positive nor a negative gradient signal. Second, every non-trivial advantage is uniformly amplified by $1/\sqrt{f_{\rm nt}}$ relative to ZPPO. When $f_{\rm nt}$ becomes small -- early in training, or when the student is too weak/too strong for most questions -- this amplification more often pushes the importance ratio outside PPO's $[1-\epsilon_{\rm low},1+\epsilon_{\rm high}]$ window and produces the visibly worse training curves reported in Sec.~\ref{sec:experiments:recipe}(ii). The role of \emph{Norm w/o Zero} is therefore not to alter any group's centered advantage $A'$, but to estimate the normalizing $\sigma$ on the sub-population $\mathcal{G}^{\star}$ that actually carries learning signal, so that the scale of the non-trivial advantages does not depend on the trivial fraction $1\!-\!f_{\rm nt}$.
\end{itemize}
The implementation (Sec.~\ref{sec:appendix:hyperparameters:zppo}, \texttt{advantage estimator} row) registers Eqs.~\eqref{eq:adv-step1}--\eqref{eq:adv-step2} with $\mathcal{S}\!=\!\mathcal{G}^{\star}$ as the ZPPO default, so all main-paper numbers use \emph{Norm w/o Zero}. The change from \emph{Norm w/ Zero} to \emph{Norm w/o Zero} is a one-line gating on $\mathrm{std}_x\!>\!0$ in the batch-statistics computation, with zero compute overhead.
